% =============================================================================
%  CONCLUSIONES
% =============================================================================
\section{Conclusiones}

La solución apunta al Nuevo Humanismo: un enfoque que busca el bien común mediante un uso responsable y ético de la tecnología. Muchos ingenieros en sistemas, programadores y diseñadores fueron los responsables de crear las plataformas que hoy causan daño a las nuevas generaciones; precisamente por eso, el cambio debe empezar desde nuestra propia formación. Como futuro ingeniero en sistemas, comprendo que la pregunta no es solo si una tecnología funciona, sino a quién beneficia y a quién perjudica su diseño. Si los pecados capitales fueron incrustados en el código por avaricia, también pueden ser desincrustados por una ética del bien común. Cambiar a los demás es difícil, pero cambiarse a uno mismo lo es aún más; por eso el primer paso es adoptar y practicar el Nuevo Humanismo de manera individual, para luego inculcarlo y generar un cambio positivo. Quizás ese sea el verdadero reto de mi generación de ingenieros: construir un mundo digital que se mueva por el bien común y no por la avaricia.

